\textbf{Triangles}
Triangle dont les côtés sont $a,b,c$, les hauteurs sont $h_a, h_b, h_c$, les médianes sont
 $m_a, m_b, m_c$, le demi-perímètre est $s = \dfrac{a + b + c}{2}$ et l'aire est $A$. $\eta=\frac{1}{4R}$\\

\begin{description}
	\item[Heron-like area formulae]
		\begin{align*} 
		A &= \sqrt{s(s-a)(s-b)(s-c)}\\
			&= \left( 4 \sqrt{\eta (\eta - h_a^{-1}) (\eta - h_b^{-1}) (\eta - h_c^{-1})} \right)^{-1} \\
			&= \dfrac{4}{3} \sqrt{s(s - m_a)(s - m_b)(s - m_c)}
		\end{align*}
	\item[Rayon du cercle circonscrit] $R = \dfrac{abc}{4A}$
	\item[Rayon du cercle inscrit] $r=\dfrac{A}{s}$
	\item[Médianes] divisent en deux sous-triangles de même aire.
		\[ m_a=\dfrac{1}{2}\sqrt{2b^2 + 2c^2 - a^2} \]
	\item[Bissectrices] divisen les angles en deux parties égales.
		\[ b_a=\sqrt{bc\left(1-\left(\dfrac{a}{b+c}\right)^2\right)}\]
	\item[Loi des sinus]
	$ \dfrac{\sin\alpha}{a}=\dfrac{\sin\beta}{b}=\dfrac{\sin\gamma}{c}=\dfrac{1}{2R}$
	\item[Loi des cosinus]
	$ a^2=b^2+c^2-2bc\cos\alpha $
	\item[Théorème de la tangente]
	$ \dfrac{a+b}{a-b}=\dfrac{\tan\dfrac{\alpha+\beta}{2}}{\tan\dfrac{\alpha-\beta}{2}} $
\end{description}

\textbf{Quadrilatères}

\textbf{Quadrilatères cycliques (inscriptibles)}
Soit un quadrilatère \textbf{cyclique} de côtés $a, b, c, d$ et diagonales $e, f$. Alors $ef = ac + bd$ et l'aire $A$ peut être calculée grâce à la

\begin{description}
	\item[Formule de Brahmagupta]
	\[A = \sqrt{(s - a)(s - b)(s - c)(s - d)}\]
	où $s = \dfrac{a + b + c + d}{2}$
\end{description}

\textbf{Coordonnées sphériques}
\[\begin{array}{cc}
x = r\sin\theta\cos\phi & r = \sqrt{x^2+y^2+z^2}\\
y = r\sin\theta\sin\phi & \theta = \texttt{acos}(z/\sqrt{x^2+y^2+z^2})\\
z = r\cos\theta & \phi = \texttt{atan2}(y,x)
\end{array}\]

\textbf{Centre de gravité:}
Soit $A$ un sous-ensemble mesurable borné de $\mathbb{R}^d$, de mesure positive. Le centre de
gravité de $A$ est le point $G = (g_1, . . . , g_d) \in \mathbb{R}^d$ dont les coordonnées sont 
définies par:
\[ 
	g_i = \frac{1}{\lambda(A)} \int_A x_i dx_1 \ldots dx_d
\]